\documentclass[11pt]{article}
\usepackage{series}

\title{Modal Triplet Theory: Foundations}
\author{Peter Nero}
\date{July 2026}

\begin{document}
\maketitle
\seriestagline

\begin{abstract}
We give a corrected functional-analytic foundation for Modal Triplet Theory
(MTT). The abstract architecture is a Hilbert bundle with three compatible
vertical structures, a joint coherent spectral projector, a stabilization
flow, and explicitly separate hypotheses for gap, invariance, existence,
contraction, truncation, and admissibility. The canonical physical realization
is a ten-dimensional bundle $M_{10}\to Y_4$ with compact Riemannian fiber
$X_6$; the central circle is bundle data and is not counted as an additional
product dimension. Strong commutation or a single total internal operator is
assumed rather than inferred from notation. Complementary-mode stability uses
a stable-semigroup estimate that remains valid for nonnormal generators.
Projected time-step fixed points are distinguished from equilibria, and
Banach, Schur--Feshbach, projector-stability, and basin-robustness statements
are given with their required domains. Stabilization time, physical time, and
renormalization scale are separated. Selection by reset is identified as a
hybrid law unless derived from continuous upper dynamics. Lorentzian signature
belongs to a hyperbolic principal symbol in a physical completion, not to a
positive Hilbert-space Gram form. A complete admissibility ledger records the
independent obligations inherited by every downstream MTT realization. A
rank-three world-in-world comparison field and the selected q79 trace-split
carrier are included as a typed geometry interface: their matching component
counts do not by themselves derive a ten-dimensional manifold, Lorentzian
spacetime, or a global intertwiner. The shared-circle claim is upgraded from
fiberwise analogy to an exact finite differential-line theorem: one universal
flat $\mathbb Z_{64}$ line pulls back coherently to the q79 SpinC determinant,
the $1+2+3$ carrier, the root-plane complex structure, and the finite Reynolds
Hessian. Boothby--Wang geometry independently identifies lens and Heisenberg
nil manifolds as parallel curved prequantum circle bundles over different
bases. These results are compatible but not identical, and neither compact
circle flow is physical Lorentzian time.
\end{abstract}

\section*{Revision note for this edition}
\addcontentsline{toc}{section}{Revision note for this edition}
\begin{description}
\item[Supersedes.] \emph{Modal Triplet Theory: Foundations}, version~7.
\item[Reason.] Version~7 correctly typed the shared circle as line-bundle
data, but it did not yet distinguish the curved Boothby--Wang Lens/Nil
realizations from the later exact flat q79 differential-line theorem, nor did
it record the resulting finite connection, holonomy, and Hessian identities.
\item[Resolution.] Version~8 retains the functional-analytic spine and adds
the universal q79 flat differential line, its SpinC/CLN/root-plane/finite-
Hessian pullbacks, the parallel Boothby--Wang realization theorem, and the
conditional polarized-section readout.
\item[Retained result.] The coherent-projector and basin-local fixed-point
architecture survives as a conditional control-and-reduction framework.
\item[Remaining boundary.] The flat finite theorem does not identify the
nonzero-Chern physical HYM connection, prove the local strain-to-q79 continuum
intertwiner, select a physical Hilbert space, or identify compact Reeb flow
with time.
\end{description}

\tableofcontents

\section{Status, scope, and logical vocabulary}

This paper defines an abstract control and reduction architecture. It does not
derive a particular quantum theory, gauge group, particle spectrum, spacetime
equation, probability law, or numerical constant. Downstream claims must use
one of the following statuses:
\begin{description}
\item[Axiom or assumption.] Input structure of a realization.
\item[Conditional theorem.] A consequence proved from listed hypotheses.
\item[Characterization.] A necessary form for objects satisfying the premises.
\item[Reconstruction or embedding.] Recovery or representation of a known
framework after compatible data are supplied.
\item[Calibration or postdiction.] Numerical agreement using target-related
input or model selection.
\item[Held-out prediction.] A quantity not used in construction, calibration,
scale choice, or branch selection.
\item[Interpretation.] A conceptual reading without theorem status.
\end{description}

The word ``physical'' is reserved for a realization equipped with a selected
state space, local evolution law, observable map, and empirical interpretation.

\section{Dimension-neutral architecture and physical realization}

\subsection{Abstract Hilbert-bundle form}

Let $Y$ be a smooth base and let $\mathcal E\to Y$ be a real or complex
separable Hilbert bundle. A state belongs to a declared Sobolev space
\[
\mathcal H=H^s(Y;\mathcal E),
\]
with $s$ chosen so that every nonlinear map and operator domain used below is
well defined. The abstract results are dimension neutral.

The modal triplet is represented by three compatible vertical structures
\[
 (\mathcal E_i,A_i,P_i),\qquad i=1,2,3,
\]
on the same internal Hilbert fiber. Here $A_i$ is a nonnegative self-adjoint
vertical operator and $P_i$ is its selected low spectral projector. The
triplet is not automatically a product of three coordinate manifolds.

\subsection{Canonical physical realization}

The canonical physical specialization is
\[
 \pi:M_{10}\longrightarrow Y_4,
 \qquad X_x=\pi^{-1}(x),qquad\dim X_x=6,
\]
where $Y_4$ is a four-dimensional globally hyperbolic Lorentzian base in the physical completion
and each compact fiber $X_x$ is Riemannian. In a local trivialization,
$M_{10}\simeq Y_4\times X_6$. Positive elliptic modal operators act vertically
on $X_6$ and do not create additional causal directions.

If $X_6$ is explicitly factorized as $F_1\times F_2\times F_3$, then
$\sum_i\dim F_i=6$. A shared phase circle is represented by a principal
$U(1)$ bundle or Hermitian line bundle $L_{\rm cen}\to X_6$. It is not appended
as a seventh independent product coordinate. Recursive nil/lens/circle
descriptions may be used as bundle or filtration data, but their dimensions
must not be added unless an actual product decomposition is proved.

\subsection{Local comparison carrier and selected q79 interface}

Let $TP$ and $TI$ be oriented Euclidean rank-three bundles over a common
base. A world-in-world comparison field is a typed section
\[
 Q_{\rm WW}\in\Gamma\!\left(\operatorname{Hom}(TP,TI)\right).
\]
After choosing local frames, $Q_{\rm WW}$ has nine matrix components. This is
not dimension multiplication: the total space of a rank-three vector bundle
over a three-dimensional base has dimension $3+3=6$.

At a nonsingular comparison background, polar decomposition separates the
infinitesimal orientation and strain directions,
\[
 \operatorname{Mat}(3,\mathbb R)
 =\mathfrak{so}(3)\oplus\operatorname{Sym}(3,\mathbb R),
 \qquad 9=3+6.
\]
Relative to a selected orthonormal flag, the strain space has the orthogonal
decomposition
\[
 \operatorname{Sym}(3,\mathbb R)
 =\mathbb RI_3\oplus\mathcal D_0\oplus\mathcal O,
 \qquad \dim(\mathbb RI_3,\mathcal D_0,\mathcal O)=(1,2,3),
\]
where $\mathcal D_0$ is traceless diagonal and $\mathcal O$ is symmetric
off-diagonal. Thus
\[
 1+3\times3=(1+3)+(1+2+3)=4+6=10
\]
is an exact component identity once one additional ordering scalar is supplied.
It is not a tangent-space decomposition of $M_{10}$, does not select a
$3+1$ Lorentzian base, and does not globalize without transition data.

For the selected degree-three q79 spectral cover, let
\[
 \mathcal A=\pi_*\mathcal O_C,
 \qquad \mathcal A=\mathcal O\oplus\mathcal A_0,
 \qquad \operatorname{rank}(\mathcal O,\mathcal A_0,\mathcal A)=(1,2,3),
\]
where $\mathcal A_0=\ker\operatorname{Tr}$. The corresponding common carrier
is
\[
 \mathcal H_{\rm CLN}
 =L_{\rm shared}\otimes(\mathcal O\oplus\mathcal A_0\oplus\mathcal A),
 \qquad \operatorname{rank}\mathcal H_{\rm CLN}=6.
\]
This proves the selected rank filtration, not its identification with the
local strain decomposition. Such an identification requires a same-source
bundle-and-connection intertwiner. The line bundle $L_{\rm shared}$ carries
common $U(1)$ phase or holonomy data; it is neither an extra product dimension
nor physical Lorentzian time.

\subsection{Universal q79 differential line and finite operator square}

The phrase ``one shared circle'' has a global meaning only after its
classifying maps and connection are retained. Let
$\mathcal L_{64}^{\rm univ}\to B_\nabla\mathbb Z_{64}$ be the universal flat
Hermitian line associated to the primitive character
\[
 \chi_1(n)=\exp(2\pi\mathrm i n/64).
\]
On the unbranched q79 sheet stack, let
$c_{\rm sheet}:B^\circ\to BS_3$ classify the sheet local system. Since
$S_3^{\rm ab}\cong\mathbb Z_2$, there is exactly one nontrivial homomorphism
to $\mathbb Z_{64}$,
\[
 h_{S_3}(\sigma)=32\,\epsilon(\sigma),
 \qquad \epsilon(\sigma)=
 \begin{cases}0,&\sigma\text{ even},\\1,&\sigma\text{ odd}.
 \end{cases}
\]
Put
\[
 c_{\rm sh}=B(h_{S_3})\circ c_{\rm sheet},
 \qquad
 (L_{\rm sh},\nabla_{\rm sh})
 =c_{\rm sh}^{*}(\mathcal L_{64}^{\rm univ},\nabla_{64}).
\]
For both admissible odd roots $r\in\{1,33\}$,
\[
 \chi_r\circ h_{S_3}=\operatorname{sgn}.
\]
Thus the associated determinant sign line and $L_{\rm sh}$ are canonically
parallel-isomorphic: their connections and every loop holonomy agree, not
merely their fiber dimensions.

\begin{theorem}[Finite shared-line and Hessian intertwiner]
\label{thm:shared-line}
On the selected q79 flat root-stack symbol, the same pullback
$L_{\rm sh}$ tensors all three trace lanes
\[
 \mathcal H_{\rm CLN}
 =L_{\rm sh}\otimes(\mathcal O\oplus\mathcal A_0\oplus\mathcal A).
\]
Its scalar holonomy commutes with the trace, trace-zero, and full projectors,
with the root-plane quarter-turn $J_{DE}$, and with the normalized finite
Reynolds operators
\[
 P_{\rm Haar}=\frac1{6}\sum_{g\in S_3}\rho(g),
 \qquad
 H_{\rm fin}=\kappa_{\rm fin}(I-P_{\rm Haar}).
\]
For the selected two-copy sheet/edge representation,
\[
 \operatorname{rank}P_{\rm Haar}=2,
 \qquad
 \operatorname{spec}(H_{\rm fin}/\kappa_{\rm fin})
 =\{0^{\times2},1^{\times4}\},
 \qquad H_{\rm TT}=\kappa_{\rm fin}I_2.
\]
The lifted Hessian is exactly
$I_{L_{\rm sh}}\otimes H_{\rm fin}$ and introduces no dimensionless fit.
\end{theorem}

\begin{proof}
The character identity gives the parallel determinant-line isomorphism.
Scalar line holonomy commutes with every finite sheet operator. Haar averaging
is an orthogonal projector, so $I-P_{\rm Haar}$ is the complementary
projector. Direct decomposition of the selected two-copy permutation
representation gives the displayed ranks and spectrum. Tensoring by the same
line preserves these identities and proves the intertwining statement.
\end{proof}

The theorem is exact at flat differential-character and finite-symbol tier.
It does not identify $\nabla_{\rm sh}$ with a nonflat physical HYM connection:
a flat trace-free carrier has vanishing real characteristic curvature, whereas
the active physical bundle has nonzero $c_2$. The continuum target is a
spectral-symbol functor and a unitary parallel Hessian comparison, not literal
equality of those connections.

\subsection{Boothby--Wang Lens and Nil realizations}

Let $(B,\omega)$ be an integral symplectic manifold. The Boothby--Wang
construction produces a principal circle bundle
\[
 U(1)\longrightarrow P_k\xrightarrow{\pi}B,
 \qquad c_1(P_k)=k[\omega],
\]
with connection/contact form $\alpha_k$ satisfying, in the chosen
normalization,
\[
 d\alpha_k=2\pi k\,\pi^*\omega.
\]
For $B=\mathbb{CP}^1$ this total space is the lens space $L(k,1)$; hence
$k=3$ gives $L(3,1)$ exactly at bundle-topology level. For an integral area
form on $T^2$, the total space is a quotient of the Heisenberg group by a
lattice and is a three-dimensional contact nilmanifold
\cite{BoothbyWang,ContactBlowup}.

Consequently Lens and Nil are parallel realizations of one prequantization
schema over different bases and curvature classes. They are not literally
nested manifolds. Nor are they one literal line bundle: the rigorous common
target is the differential classifier $BU(1)_\nabla$ with separate maps
\[
 c_{\rm Lens}:\mathbb{CP}^1\to BU(1)_\nabla,
 \qquad
 c_{\rm Nil}:T^2\to BU(1)_\nabla.
\]
The flat q79 classifying map factors through
$B_\nabla\mathbb Z_{64}\to BU(1)_\nabla$. Because the Boothby--Wang
connections above have nonzero curvature while the q79 root-stack line is
flat on $B^\circ$, the constructions cannot be identified connection by
connection. A stronger MTT ``same circle'' theorem would have to select the
three classifying maps and coherent comparison $2$-cells on a declared
correspondence space.

After a polarization, weight-$m$ equivariant functions on $P_k$ correspond to
sections of the associated line $L^m$. In a positive holomorphic polarization
this gives the familiar readout
\[
 \mathcal H_m=H^0(B,L^m).
\]
This is a conditional geometric-quantization Hilbert space, not a Hilbert
space selected by the abstract MTT axioms. The periodic Reeb flow is the
vertical phase action and must not be identified with noncompact physical
time \cite{Woodhouse}.

\section{Three independent evolution parameters}

The foundation distinguishes:
\begin{enumerate}
\item a stabilization flow $R_\tau$ with control parameter $\tau\ge0$;
\item a physical propagator $U(t_2,t_1)$, supplied only by a selected physical
completion; and
\item a renormalization or coarse-graining scale $\mu$.
\end{enumerate}
No equality among $\tau$, $t$, and $\log\mu$ is assumed. A theorem relating any
two of them must specify the map, units, domain, and approximation error.

We write the stabilization equation as
\[
 \partial_\tau\Psi=F(\Psi),\qquad R_\tau(\Psi_0)=\Psi(\tau).
\]
Well-posedness is imposed on a declared interval and invariant domain; global
well-posedness is not part of the abstract foundation unless proved in a
specific realization.

\section{Joint vertical operator and coherent projector}

\begin{assumption}[Joint spectral structure]\label{ass:joint}
For each base point, the $A_i$ are nonnegative self-adjoint operators whose
spectral measures strongly commute. Their quadratic forms have a common dense
domain. Equivalently, a realization may provide one nonnegative self-adjoint
total internal operator $A_{\rm int}$ directly.
\end{assumption}

Under strong commutation define
\[
 P=\prod_{i=1}^3\mathbf1_{I_i}(A_i),\qquad Q=I-P,
\]
for declared isolated low spectral sets $I_i$. The product is then an
orthogonal projector independent of ordering. Base-only coefficients or
separate variable names do not, by themselves, prove strong commutation.

For harmonic selection one may instead use the form sum
$A_{\rm int}=A_1+A_2+A_3$. Nonnegativity gives
\[
 \ker A_{\rm int}=\bigcap_i\ker A_i,
\]
provided the form sum is well defined. The coherent projector is then the
spectral projector of $A_{\rm int}$ at zero.

\begin{assumption}[Internal gap and Sobolev boundedness]\label{ass:gap}
There is $\lambda_{\rm int}>0$ such that, in quadratic-form sense,
\[
 A_{\rm int}\succeq\lambda_{\rm int}Q,
\]
and $P,Q$ extend boundedly to every Sobolev space used by the dynamics.
\end{assumption}

The internal gap separates the selected fiber cluster from complementary
fiber modes. It does not by itself imply invariance under $R_\tau$, existence
of a fixed point, coherent contraction, suppression of arbitrarily high
four-dimensional energy, or selection of a physical state.

\section{Stable complementary dynamics}

Let $\Psi_\ast$ be a reference state at which the stabilization vector field is
Fr\'echet differentiable, and set $L=DF(\Psi_\ast)$. The stable sign convention
is that $L_{QQ}=QLQ$ generates decay:
\[
 \|e^{\tau L_{QQ}}\|\le M_Qe^{-\omega_Q\tau},
 \qquad M_Q\ge1,quad\omega_Q>0.
 \label{eq:q-semigroup}
\]
This estimate, rather than spectral abscissa alone, is authoritative for a
nonnormal generator.

\begin{proposition}[Gap-to-decay under a bounded perturbation]
Suppose $Q\mathcal H$ is invariant and
\[
 L_{QQ}=-\kappa A_{\rm int}|_{Q\mathcal H}+B_Q,
 \qquad \kappa>0,quad B_Q\in\mathcal B(Q\mathcal H).
\]
If $\omega_Q:=\kappa\lambda_{\rm int}-\|B_Q\|>0$, then
\[
 \|e^{\tau L_{QQ}}\|\le e^{-\omega_Q\tau}.
\]
\end{proposition}

\begin{proof}
The self-adjoint semigroup generated by
$-\kappa A_{\rm int}|_{Q\mathcal H}$ has norm at most
$e^{-\kappa\lambda_{\rm int}\tau}$. The bounded perturbation estimate gives
the additional factor $e^{\|B_Q\|\tau}$.
\end{proof}

For a general sectorial or nonnormal $L_{QQ}$, the constants $M_Q$ and
$\omega_Q$ must be proved directly. A positive vertical operator cannot be
identified with the stabilization generator without the minus sign and the
lower-order terms.

\section{Independent logical gates}

The following gates are independent and must not be collapsed:
\begin{description}
\item[Gap.] Spectral separation for $A_{\rm int}$.
\item[Projector stability.] Persistence and regularity of $P$ under parameter
or curvature variation.
\item[Invariance.] $R_\tau(P\mathcal H)\subseteq P\mathcal H$, equivalently
$QF(Pu)=0$ in a differentiable autonomous realization.
\item[Existence.] A fixed-point theorem applies on a declared invariant domain.
\item[Equilibrium identification.] A projected time-step fixed point is shown
to be stationary.
\item[Contraction.] The coherent reduced map is strictly contractive.
\item[Truncation.] The influence of $Q$ on $P$ is quantitatively bounded.
\item[Selection.] A continuation or reset rule is supplied at admissibility
exit.
\end{description}

No gate in this list follows solely from the gate preceding it.

\section{Projected fixed points and equilibria}

Let $K\subset P\mathcal H$ be nonempty, closed, bounded, and convex and define
\[
 T_\tau=P R_\tau|_K.
\]

\begin{theorem}[Projected time-step existence]\label{thm:existence}
Assume $T_\tau(K)\subset K$ and either:
\begin{enumerate}
\item $T_\tau:K\to K$ is continuous and compact; or
\item $T_\tau$ is continuous and condensing for a declared measure of
noncompactness.
\end{enumerate}
Then there is $u_\ast\in K$ with $T_\tau u_\ast=u_\ast$.
\end{theorem}

\begin{proof}
Apply Schauder in the first case and Darbo--Sadovskii in the second.
\end{proof}

The conclusion is a \emph{projected time-step fixed point}. It is not named an
equilibrium merely because $P R_\tau u_\ast=u_\ast$.

\begin{theorem}[Strict-Lyapunov equilibrium promotion]\label{thm:promotion}
Suppose $P\mathcal H$ is invariant, so that
$P R_\tau u_\ast=R_\tau u_\ast=u_\ast$, and along the orbit
\[
 \mathcal C(R_\tau u)-\mathcal C(u)
 =-\int_0^\tau\mathcal D(R_su)\,ds,
 \qquad\mathcal D\ge0,
\]
where $\mathcal D(v)=0$ exactly when $F(v)=0$. Then the fixed point from
Theorem~\ref{thm:existence} is an equilibrium.
\end{theorem}

\begin{proof}
Endpoint recurrence makes the left side zero. Nonnegativity and continuity
force $\mathcal D=0$ along the orbit; strictness gives $F(u_\ast)=0$.
\end{proof}

\begin{theorem}[Banach gate]\label{thm:banach}
If $K$ is complete, $T_\tau(K)\subset K$, and
\[
 \|T_\tau u-T_\tau v\|\le q\|u-v\|,
 \qquad0\le q<1,
\]
then $T_\tau$ has a unique fixed point in $K$, and its iterates converge to
that point. The theorem does not apply without the invariant complete domain.
\end{theorem}

\section{Schur--Feshbach reduction with domains}

Let $\mathcal H=P\mathcal H\oplus Q\mathcal H$. Consider a closed block
operator
\[
 L=\begin{pmatrix}L_{PP}&L_{PQ}\\L_{QP}&L_{QQ}\end{pmatrix}
\]
with domain $P\mathcal H\oplus\mathcal D(L_{QQ})$. Assume:
\begin{enumerate}
\item $0\in\rho(L_{QQ})$ and
$L_{QQ}^{-1}:Q\mathcal H\to\mathcal D(L_{QQ})$ is bounded in the graph norm;
\item $L_{PP}$ and $L_{QP}:P\mathcal H\to Q\mathcal H$ are bounded; and
\item $L_{PQ}:\mathcal D(L_{QQ})\to P\mathcal H$ is graph-norm bounded, so
$L_{PQ}L_{QQ}^{-1}$ is bounded.
\end{enumerate}

\begin{theorem}[Schur--Feshbach equation]\label{thm:feshbach}
Under these hypotheses, solving $L(p,q)=(f_P,f_Q)$ is equivalent to
\[
 S p=f_P-L_{PQ}L_{QQ}^{-1}f_Q,
 \qquad
 q=L_{QQ}^{-1}(f_Q-L_{QP}p),
\]
where
\[
 S=L_{PP}-L_{PQ}L_{QQ}^{-1}L_{QP}.
\]
If $\|L_{QQ}^{-1}\|\le C_Q/\omega_Q$, then
\[
 \|S-L_{PP}\|
 \le\|L_{PQ}L_{QQ}^{-1}\|\,\|L_{QP}\|
 \le\frac{C_Q\|L_{PQ}\|_{\rm gr}\|L_{QP}\|}{\omega_Q}.
\]
\end{theorem}

\begin{proof}
The $Q$ block equation gives the displayed formula for $q$. Substitution in
the $P$ block equation gives $S$. The bound follows from the declared graph
norm and inverse estimates.
\end{proof}

This is a local linear reduction near the reference state. Nonlinear
truncation additionally requires control of nonlinear remainders and of the
time interval on which the eliminated sector remains small.

\section{Projector stability and basin-local robustness}

\begin{proposition}[Riesz-projector stability]
Let $A(\epsilon)$ be a norm-resolvent-continuous family and let a contour
$\Gamma$ remain in the resolvent set while enclosing one isolated cluster.
Then
\[
 P(\epsilon)=\frac{1}{2\pi i}\oint_\Gamma(z-A(\epsilon))^{-1}\,dz
\]
is norm continuous and has constant finite rank. This proves projector
stability, not dynamical stability of a fixed point.
\end{proposition}

\begin{theorem}[Basin-local fixed-point robustness]\label{thm:robust}
Let $T,\widetilde T:K\to K$ be contractions on the same complete invariant
domain with contraction constant at most $q<1$. If
\[
 \sup_{u\in K}\|T(u)-\widetilde T(u)\|\le\varepsilon,
\]
then their fixed points satisfy
\[
 \|u_\ast-\widetilde u_\ast\|\le\frac{\varepsilon}{1-q}.
\]
\end{theorem}

\begin{proof}
Insert and subtract $T(\widetilde u_\ast)$ and absorb the resulting
$q\|u_\ast-\widetilde u_\ast\|$ term.
\end{proof}

This is the justified foundation for a local robustness or universality
claim. It is not global universality across arbitrary microscopic models.

\section{Projection, descent, and recovery types}

Let $r:X\to Y_{\rm eff}$ be a surjective reduction map on a microscopic state
space $X$, and let $\Phi:X\to X$ be a microscopic step.

\begin{theorem}[Autonomous descent criterion]\label{thm:descent}
There exists a unique reduced map $\overline\Phi:Y_{\rm eff}\to Y_{\rm eff}$
satisfying
\[
 \overline\Phi\circ r=r\circ\Phi
\]
if and only if
\[
 r(x)=r(x')\quad\Longrightarrow\quad r(\Phi x)=r(\Phi x').
\]
\end{theorem}

\begin{proof}
Necessity follows by applying $\overline\Phi$. For sufficiency define
$\overline\Phi(r(x))=r(\Phi x)$; the implication makes this independent of the
representative, and surjectivity gives uniqueness.
\end{proof}

A section $s:Y_{\rm eff}\to X$ obeys $r\circ s=\operatorname{id}$ and chooses
one representative. It does not recover every microscopic state. Exact
microscopic recovery would require $s\circ r=\operatorname{id}_X$, which is
possible only when $r$ is injective. An effective merger combines reduced
descriptions and is neither kind of inverse unless separately typed and
proved.

If $r$ is fiberwise over the base and the microscopic generator is local in
base variables, a reduced local generator may descend when the invariance
criterion holds and all coefficients depend locally and smoothly on the base
jet. A projection nonlocal in base variables does not inherit locality merely
from being a projector.

\section{Admissibility and hybrid selection}

Let $m_j(u)$ be declared continuous margins and define
\[
 \mathcal A_\varepsilon
 =\{u:m_j(u)\ge\varepsilon\text{ for all }j\},
 \qquad
 \mathfrak B(u)=\max_j[-m_j(u)].
\]
Then $u\in\mathcal A_0$ exactly when $\mathfrak B(u)\le0$. This is a diagnostic
encoding of declared constraints, not automatically an energy or force.

At a first exit from $\mathcal A_0$, three logically different constructions
are possible:
\begin{enumerate}
\item stop the reduced description;
\item continue the upper dynamics and derive a new reduced chart; or
\item impose a reset $S:\partial\mathcal A_0\to\mathcal A_0$.
\end{enumerate}
The third option defines a hybrid dynamical system. It changes the continuation
law and must separately prove existence, measurability, conservation, and any
probability assigned to alternative reset outcomes. Calling the reset
``selection'' does not make it part of the original flow.

\section{Complete admissibility-margin ledger}

Every realization must record the following entries independently:
\begin{enumerate}
\item \textbf{Geometry/domain:} base, fiber, dimensions, bundle regularity,
operator domains, and boundary conditions.
\item \textbf{Joint spectral structure:} strong commutation or a selected
total operator, the spectral contour, rank, and Sobolev boundedness of $P$.
\item \textbf{Internal gap:} $\lambda_{\rm int}$ and its uniformity domain.
\item \textbf{Complementary stability:} $M_Q,\omega_Q$ and the generator sign.
\item \textbf{Leakage/invariance:} $QF(Pu)$ or a quantitative leakage bound.
\item \textbf{Existence gate:} the invariant set $K$ and the compact,
condensing, or contraction theorem actually used.
\item \textbf{Equilibrium gate:} stationarity equation or strict Lyapunov
promotion.
\item \textbf{Coherent contraction:} norm, basin, and factor $q<1$.
\item \textbf{Truncation:} mixing blocks, inverse/domain bounds, nonlinear
remainder, and error interval.
\item \textbf{Projector robustness:} perturbation topology and resolvent
contour.
\item \textbf{Admissibility:} each margin $m_j$, its units, and its source.
\item \textbf{Continuation/selection:} stop, upper continuation, or hybrid
reset, including conservation and probability data.
\item \textbf{Physical hyperbolicity:} principal symbol, constraints, causal
domain, and relation between $t$ and $\tau$.
\item \textbf{Scale separation:} record the internal gap
$\lambda_{\rm int}$, coherent contraction scale, four-dimensional cutoff,
curvature scale, and RG scale $\mu$.
\item \textbf{Numerical provenance:} source-independent inputs, fitted inputs,
branch selection, uncertainty, and held-out outputs.
\end{enumerate}

No single scalar ``coherence scale'' may replace this ledger unless a theorem
derives all identifications with dimensions and errors.

\section{Lorentzian physical completion}

The Hilbert inner product and every Gram tensor constructed from it are
positive semidefinite. They cannot acquire Lorentzian signature. In a local
physical completion, causal signature is instead read from the principal
symbol. For a second-order field equation this has the schematic form
\[
 \sigma_{\rm pr}(x,\xi)=g^{\mu\nu}(x)\xi_\mu\xi_\nu I+\text{gauge blocks}.
\]
Lorentzian hyperbolicity, constraint propagation, and a domain-of-dependence
theorem must be verified for the selected equations. The canonical base
dimension $3+1$ is part of the canonical FP physical realization; it is not
derived by the abstract Hilbert-bundle theory. In particular, the component
identity $1+3\times3=4+6$ neither selects a four-dimensional causal base nor
determines its signature.

\section{Foundation synthesis theorem}

\begin{theorem}[Scoped MTT foundation]\label{thm:foundation}
Assume the dimension-neutral Hilbert-bundle architecture, joint spectral
structure, internal gap, stable complementary semigroup, and the separately
listed hypotheses for a chosen fixed-point theorem. Then the model possesses
a bounded coherent decomposition and a projected time-step fixed point. Under
strict Lyapunov and invariance hypotheses that point is an equilibrium. Under
the Banach gate it is unique in the declared basin. Under the Schur--Feshbach
hypotheses the local linear complementary sector can be eliminated with the
displayed error bound. Under contraction and map-closeness hypotheses the
fixed point is basin-locally robust. Autonomous reduced dynamics exists
exactly under the descent criterion.
\end{theorem}

None of these conclusions supplies a physical probability law, Lorentzian
field equation, quantum representation, particle spectrum, Standard Model
matching, cosmology, or numerical prediction. Those are downstream
obligations governed by the admissibility ledger. Nor does equality of the
local $1+2+3$ strain dimensions and the q79 $1+2+3$ trace-split ranks prove the
same-source intertwiner needed to identify their metrics, connections, and
vertical operators. Theorem~\ref{thm:shared-line} closes the common flat line
and finite Hessian square inside the q79 carrier; it does not close that
local-to-continuum identification.

\section{Conclusion}

The corrected Foundation provides a typed and noncircular spine for the MTT
corpus. It identifies what the triplet and projector mean, how complementary
stability is proved, which fixed-point conclusion is available, how reduction
is typed, and where physical assumptions enter. Its value is precisely this
separation: later realizations can now be tested against explicit gates rather
than inheriting conclusions from an undifferentiated appeal to coherence.
The shared-circle sector now illustrates the intended discipline particularly
well: a curved Boothby--Wang schema, an exact flat q79 differential line, and
a conditional polarized Hilbert readout coexist without being conflated.

\begin{thebibliography}{99}

\bibitem{FixedPoints}
P.~Nero,
\newblock \emph{Fixed Points I--VI}, corrected theorem spine,
\newblock revised editions, 2026.

\bibitem{Kato}
T.~Kato,
\newblock \emph{Perturbation Theory for Linear Operators},
\newblock Springer, 1995.

\bibitem{EngelNagel}
K.-J.~Engel and R.~Nagel,
\newblock \emph{One-Parameter Semigroups for Linear Evolution Equations},
\newblock Springer, 2000.

\bibitem{Deimling}
K.~Deimling,
\newblock \emph{Nonlinear Functional Analysis},
\newblock Springer, 1985.

\bibitem{BoothbyWang}
W.~M.~Boothby and H.~C.~Wang,
\newblock On contact manifolds,
\newblock \emph{Annals of Mathematics} 68 (1958), 721--734.

\bibitem{ContactBlowup}
R.~Casals, D.~M.~Pancholi, and F.~Presas,
\newblock Contact blow-up,
\newblock \emph{Expositiones Mathematicae} 33 (2015), 97--123.

\bibitem{Woodhouse}
N.~M.~J.~Woodhouse,
\newblock \emph{Geometric Quantization}, second edition,
\newblock Oxford University Press, 1992.

\bibitem{SharedLinePacket}
P.~Nero,
\newblock \emph{q79 Universal Shared Differential Line and Finite-Operator
Intertwiner}, executable theorem packet, 2026.

\end{thebibliography}

\end{document}
